\section{Introduction}
\label{sec:introduction}
% NOTE (this revision pass): rewritten for shorter sentences, explicit
% transitions between paragraphs, and one new comparison table (Table
% \ref{tab:related-approaches}) that replaces the old wall of "(cite) ..."
% bullets with a scannable summary. Content is unchanged from the previous
% version -- this is a style/structure pass, not a new round of fact-finding.
% Each subsection still opens with an italic quote block reproducing your
% original main.tex bullet points verbatim.

\subsection{Background}
\begin{quote}\small\itshape
\begin{itemize}
\item EHR Foundation Model (FM) arises nowaday, with potential to be assistive clinical decision tool
\item they often use standarized clinical vocabularies such as ICD, RXNORM, SNOMED, LONIC
\item standarized clinical vocabularies /ontologies DAG, are used to normalize clinical naming, and improve the interoperability
\item SNOMED CT is a DAG, comprehensive 400k + concepts (covering very broad clinical domains, diorders, medication, morphologies etc...) and semantically rich (composed by 100+ relationship types, have their own way to compose machine readable semantic expression), mappable to other clincial terminologies.
\end{itemize}
\end{quote}
Electronic health records (EHRs) are the data generated from a patient's routine healthcare encounters. They capture diagnoses, prescriptions, surgical procedures, and other clinical events. Hospitals produce large volumes of EHR data daily, and with the growth of Artificial Intelligence (AI), this data has gained increasing value for training clinical decision support systems.

Foundation models (FMs) are a recent class of AI models that learn from large amounts of patient history, with the goal of eventually assisting clinical decisions \citeneeded{self-spervised scoping,li2020behrt,yang2023transformehr}

Much of the EHR data used to train FMs are clinical events, encoded using standardized medical vocabularies. Several such vocabularies are in wide use, including RxNorm, LOINC, SNOMED~CT, and ICD \citeneeded{standardized clinical vocabulary / interoperability reference}. They not only standardize clinical terminology but also enable interoperability across healthcare institutions.

Several of these vocabularies are more than flat lists of codes: they are ontologies, organized as directed acyclic graphs (DAGs), in which each concept is connected to other concepts that define it. SNOMED CT is one such ontology, and currently the most comprehensive clinical terminology in use. It is maintained on a continuous basis, with a new International Edition released monthly, and contains more than 400{,}000 concepts, covering disorders, procedures, medications, and more. Its concepts are connected by over 100 relationship types and can be mapped onto most other major clinical terminologies \citeneeded{SNOMED International concept model / release statistics}.

\subsection{Gap}
\begin{quote}\small\itshape
\begin{itemize}
\item training FM needs either large amount of patient data (but patient data is difficult to get)
\item or need adequate size of clinical vocabularies. (but medical vocabularies are large size)
\item Current methods are:
\item (cite) LLM byte-pair on categorical medical concept: data-driven, no semantic preservation, highly dependant to data quality and cpncept orders.
\item (cite) keep highest prevalent top k concepts, highest entropy contribution top k concepts: data-driven, loss of considerable clinical information.
\item (cite) troncate/reduce leaf concepts to their parents concepts: data-driven, loss of medical granularity.
\item (cite) or just keep all concepts as it is, if we have lots of patient data. but since concept usage distribution are usually long tail, many rare concepts can not be learnt in a proper way.
\item (cite) learnable embedding space codebook: loss of controillability and interpretability
\item all of above mentioned methods are ignoring the semantic richness of clinical vocabulary DAG when try to deal with vocabulary size.
\end{itemize}
\end{quote}
One of the major difficulties hindering EHR-pretrained FMs is the trade-off between the size of the clinical vocabulary and the amount of data available to learn a good representation for each concept in it. This trade-off follows directly from how these models are built. Modern EHR-trained FMs are almost all transformer-based sequence models that operate over tokens via self-attention \citeneeded{vaswani2017attention}, where each clinical concept is assigned its own learnable embedding that is refined only through the concept's own occurrences during training. Two concepts that are clinically similar but distinct in the vocabulary therefore start out, and often remain, unrelated in the model's representation space, since nothing in a standard transformer assumes they should behave alike unless the training data itself supplies enough evidence for the model to learn that they do, and generating that evidence requires a large amount of patient data, which is expensive and often hard to access. This mirrors the vocabulary-size/data-budget scaling relationship recently characterized for language models, where the vocabulary size that best matches a given training budget cannot be chosen independently of the amount of available data \citeneeded{tao2024scaling}. However, reducing the vocabulary to a size appropriate to the available training data is not trivial. Concept usage in real EHR data follows a long-tailed distribution \citeneeded{zhang2023longtailed}, so most concepts remain too rare to learn well regardless of how the vocabulary is chosen.

Several strategies already exist for shrinking the vocabulary, and Table~\ref{tab:related-approaches} summarizes them side by side. Byte-pair and other subword encodings merge concepts by co-occurrence statistics alone \citeneeded{sennrich2016bpe, kudo2018sentencepiece}: they need no ontology, but the resulting vocabulary depends heavily on data quality and on the arbitrary order in which concepts appear. Keeping only the most frequent, or highest-information, concepts (top-$k$ pruning) drops a significant amount of clinically relevant information simply because it is rare \citeneeded{prevalence/entropy-based vocabulary pruning in clinical ML}. Rolling leaf concepts up to their ontology parent shrinks the vocabulary but loses clinical detail: a specific diagnosis collapses into a broader one \citeneeded{ontology roll-up / ICD truncation approach}. Keeping the full vocabulary avoids all of these losses, but only postpones the long-tail problem to the model, and only works if enough data is available \citeneeded{EHRSHOT?}. Finally, learnable embedding-space codebooks compress the vocabulary without any hand-written rules \citeneeded{su2025medtok, vandenoord2017vqvae}, but the resulting tokens are then just indices into a learned space: they no longer correspond to concepts a person can recognize or check.

\begin{table}[h]
\centering
\small
\begin{tabular}{p{2.7cm}p{1.8cm}p{2.0cm}p{4.5cm}}
\toprule
\textbf{Approach} & \textbf{Use of semantic} & \textbf{Controllability} & \textbf{Main limitation} \\
\midrule
Byte-pair / subword encoding & No & No (merges are statistical, not semantic) & Sensitive to data quality and concept order \\
Frequency / entropy top-$k$ & No & Yes (kept items are still real concepts) & Drops clinically relevant but rare concepts \\
Hierarchy truncation / roll-up & Partially (parent--child links only) & Yes, but coarser & Loses clinical granularity \\
Full vocabulary (no reduction) & N/A & Yes & Long tail is never solved, just handed to the model \\
Ontology-guided attention (e.g.\ GRAM) & Yes (via ancestor attention) & Partial (representation is learned, not a fixed token) & Improves rare-code learning but does not fix vocabulary size itself \\
Learnable codebook (e.g.\ MedTok) & Yes (e.g.\ via a graph encoder) & No (opaque codebook indices) & Not human-auditable \\
\textbf{Ours} & \textbf{Yes, directly} & \textbf{Yes, fully traceable} & (this paper's contribution) \\
\bottomrule
\end{tabular}
\caption{How existing vocabulary-reduction and representation-learning strategies compare on the two properties this paper cares about most: whether they use the ontology's graph structure, and whether the resulting tokens can be inspected by a person.}
\label{tab:related-approaches}
\end{table}

Table~\ref{tab:related-approaches} makes the gap concrete: no existing approach combines both properties to produce a token set that is both semantically motivated and inspectable/controllable.

\subsection{Our method and contribution}
\begin{quote}\small\itshape
\begin{itemize}
\item applicable to DAG knowledge graph.
\item tokens selection: Represent concepts by smaller vocabulary size using greedy token selection approach
\item way to tokenize: Represent concepts by a tree/subgraph preserving the structure of the concepts the semantic value.
\item By apply to SNOMED CT concepts used in our local institut real data, we show our method can control how we tokenize and with what we tokenize (controllability and interpretability)
\item (maybe) if any mapping inter-ontology exist, we can use concepts in 1 ontology to represent/tokenize concepts in the other ontologies (interoperability)
\end{itemize}
\end{quote}
Building on this gap, we propose a method for selecting and building tokens directly on a concept DAG such as SNOMED~CT. It has two parts. First, we choose a small set of tokens with a greedy algorithm that maximizes how well they cover the meaning of the target concepts (Section~\ref{sec:method}). This lets us fix the vocabulary size in advance, instead of it falling out of a training run. Second, instead of representing a concept as an unordered set of the tokens it maps to, we represent it as a \emph{context tree}: a small subgraph that records which relationship led to each token.

We apply this method to the SNOMED CT concepts mapped by our institution's terminology team. The result, unlike every approach in Table~\ref{tab:related-approaches}, is directly controllable and inspectable: for any concept, a domain expert can see exactly which relationships were used to reach which token. Such a tokenizing method makes FM training more auditable.